% !TEX program = xelatex
% Working technical rider; venue specifications remain to be confirmed.
\documentclass[11pt,letterpaper]{article}
\usepackage{culturehub-packet}
\renewcommand{\packetlabel}{TECHNICAL RIDER}
\hypersetup{pdftitle={AC Presents 2 New Pieces — Technical Rider}}

% Logotype title for a work, set like a game title.
\newcommand{\worktype}[3]{%
  {\aclight\fontsize{11pt}{13pt}\selectfont\color{acpink}#1\par}%
  \vspace{0.12em}%
  {\acbold\fontsize{37pt}{38pt}\selectfont\color{acdark}#2\par}%
  \vspace{0.16em}%
  {\aclight\fontsize{14pt}{17pt}\selectfont\color{acpurple}#3\par}%
  \vspace{0.30em}\color{acpink}\rule{\linewidth}{1.15pt}\color{black}\vspace{0.10em}%
}
% Header bar for a specification block.
\newcommand{\specbar}[1]{%
  \vspace{0.35em}%
  \begingroup\setlength{\fboxsep}{5pt}%
  \colorbox{acdark}{\color{white}\acbold\fontsize{11pt}{13pt}\selectfont #1}%
  \endgroup\vspace{0.15em}\par}
\newcommand{\spectable}{\fontsize{9.6pt}{12.2pt}\selectfont}

\begin{document}
\packettitle{AC Presents 2 New Pieces}{Technical rider · CultureHub LA · September 16 to 25, 2026}
\statushold{SEPT 24 CONFIRMED · SEPT 19 PENDING}

\vspace{0.45em}
\pageclaim{Two pieces to install in the room. Each has its own machines, signal plan, and system requirements.}

\tighttable
\begin{tabularx}{\linewidth}{L{0.28\linewidth}L{0.28\linewidth}Y}
\toprule
\textbf{Date} & \textbf{Event} & \textbf{Shape} \\
\midrule
Sat Sep 19, 2:00–3:30 p.m. & Grokaesthetic Workshop & Participatory; bring your own computer \\
Thu Sep 24, 7:00 p.m. & Piece 1 and Piece 2 & Performance, conversation, free play \\
\bottomrule
\end{tabularx}
\normalsize

\vspace{0.3em}
Thursday September 24 is confirmed. Saturday September 19 is still open: the
guest set has been cut, and with it some of the case for two nights. This rider
covers both; if they merge, the workshop's requirements move to Thursday
afternoon and nothing else changes.

\section{What gets installed}
\tighttable
\begin{tabularx}{\linewidth}{L{0.24\linewidth}L{0.26\linewidth}Y}
\toprule
\textbf{System} & \textbf{Machines} & \textbf{What the house provides} \\
\midrule
Workshop & Participant laptops and phones, plus loaners & Guest Wi-Fi for 25–35 devices with downloads permitted; participant power; stereo minimum \\
\textbf{Piece 1} \newline \textit{Note(s)pat(ial) Native} & Six surplus laptops booted from AC OS USB drives & Six discrete inputs to the Kalio system; projection \\
\textbf{Piece 2} \newline \textit{The MacNeoPolitan Trio} & Three MacBook Neos running Menu Band on macOS & Three discrete full-range positions; outbound NTP \\
\bottomrule
\end{tabularx}
\normalsize

\clearpage
\worktype{Event one · Saturday}{Grokaesthetic Workshop}{Bring your own computer. Leave playing one.}

\begin{center}
\begin{tikzpicture}[
  every node/.style={font=\sffamily\bfseries\fontsize{9.5pt}{11pt}\selectfont,align=center},
  way/.style={draw=accyan,very thick,rounded corners=3pt,fill=accyan!8,minimum width=40mm,minimum height=14mm},
  need/.style={draw=acpink,very thick,rounded corners=3pt,fill=acpink!8,minimum width=40mm,minimum height=12mm},
  room/.style={draw=acpurple,very thick,rounded corners=3pt,fill=acpale,minimum width=74mm,minimum height=14mm},
  flow/.style={-{Stealth[length=3mm]},very thick,draw=acdark}
]
\node[way] (a) at (-4.75,0) {browser\\any device};
\node[way] (b) at (0,0) {Menu Band\\a Mac};
\node[way] (c) at (4.75,0) {AC OS\\USB stick};
\node[need] (an) at (-4.75,-1.95) {open internet};
\node[need] (bn) at (0,-1.95) {app downloads\\permitted};
\node[need] (cn) at (4.75,-1.95) {permission to\\boot from USB};
\node[room] (r) at (0,-4.0) {the room plays together\\house surround system};
\draw[flow] (a) -- (an);
\draw[flow] (b) -- (bn);
\draw[flow] (c) -- (cn);
\draw[flow] (an) -- (r);
\draw[flow] (bn) -- (r);
\draw[flow] (cn) -- (r);
\end{tikzpicture}
\end{center}

Ninety minutes, hands-on, free and open to the public. Participants arrive with
their own laptops and phones, get Aesthetic.Computer running on them, and the
room finishes by playing together as one ensemble. The USB route installs
nothing on a participant's disk and alters nothing.

\specbar{SYSTEM REQUIREMENTS}
\spectable
\begin{tabularx}{\linewidth}{L{0.20\linewidth}Y}
\toprule
\textbf{Network} & Guest Wi-Fi for 25–35 simultaneous devices \textbf{with app downloads permitted}. The highest-risk item here: a captive portal or a download block turns the middle route off at the door. \\
\textbf{Permissions} & Participant laptops may be booted from USB in the space. \\
\textbf{Power} & Strips reaching the seating, enough for about 20 laptops. Assume nobody arrives charged. \\
\textbf{Audio} & Stereo playback minimum; the surround system preferred. Six-channel routing is welcome but not required. \\
\textbf{Room} & Tables or floor seating arranged so people can see each other rather than a stage, plus positions for the artist's loaner machines. A microphone, or a room quiet enough to teach in. \\
\textbf{Catering} & A snack table with a nearby bin. A real requirement: the workshop is part hangout by design. \\
\bottomrule
\end{tabularx}
\normalsize

\section{Access}
No programming or musical experience is assumed, and loaner machines mean no
one is turned away for lack of hardware. Please advise on accessible seating
and whether the space suits participants who cannot sit on the floor.

\clearpage
\worktype{Piece 1 · Thursday}{Note(s)pat(ial) Native}{Six laptops carry sound around the room while it turns.}

\begin{center}
\begin{tikzpicture}[
  node distance=8mm and 8mm,
  every node/.style={font=\sffamily\bfseries\fontsize{10pt}{11pt}\selectfont,align=center},
  box/.style={draw=acpurple,very thick,rounded corners=3pt,fill=acpale,minimum width=34mm,minimum height=14mm},
  laptop/.style={draw=accyan,very thick,rounded corners=2pt,fill=accyan!8,minimum width=38mm,minimum height=14mm},
  house/.style={draw=acpink,very thick,rounded corners=3pt,fill=acpink!8,minimum width=39mm,minimum height=15mm},
  flow/.style={-{Stealth[length=3mm]},very thick,draw=acdark}
]
\node[box] (control) {live control\\rotation · speed\\trajectory · routing};
\node[box,right=8mm of control] (network) {local network\\clock · transport\\spatial state};
\node[laptop,right=8mm of network] (laptops) {six laptops\\1 · 2 · 3 · 4 · 5 · 6};
\node[house,below=13mm of laptops] (house) {house spatial path\\verified 5.1 topology};
\node[house,left=8mm of house] (audience) {audience inside\\the rotating field};
\draw[flow] (control) -- (network);
\draw[flow] (network) -- (laptops);
\draw[flow] (laptops) -- node[right,font=\sffamily\bfseries\fontsize{9pt}{10pt}\selectfont] {six feeds} (house);
\draw[flow] (house) -- (audience);
\end{tikzpicture}
\end{center}

Six surplus grade-school laptops, each booted from an AC OS USB drive, carry
independent sound bodies around the audience. The artist drives rotation,
trajectory, speed, and voice routing live from a control laptop, and the
spatial score is projected while it plays.

\specbar{SYSTEM REQUIREMENTS}
\spectable
\begin{tabularx}{\linewidth}{L{0.20\linewidth}Y}
\toprule
\textbf{Machines} & Six surplus laptops booted from AC OS USB drives, plus one control laptop. Artist-supplied, with two spares. \\
\textbf{Audio} & \textbf{Six discrete inputs} to the Kalio 5.1 system, from an interface providing six line-level outputs, with DI boxes, cabling, and adapters to match. \\
\textbf{Sixth channel} & Bandwidth confirmed. If it is LFE-only, the field uses the five full-range channels and sends only derived low-frequency material to the LFE, or takes an additional full-range speaker if the venue supplies and approves one. \\
\textbf{Network} & Dedicated router or switch, isolated from public traffic, carrying clock, transport, and spatial control. \\
\textbf{Power} & Seven positions: six laptops and the control laptop. \\
\textbf{Space} & Six low tables or plinths arranged around the audience field, with the audience inside it. \\
\textbf{Video} & Projector, screen or projection wall, and an HDMI input near the performance position, for the live spatial score. \\
\textbf{Will not run with} & Routing hard-coded as six stationary voices. The piece performs spatial rotation by moving sound bodies across the available full-range outputs. \\
\bottomrule
\end{tabularx}
\normalsize

\specbar{SIGNAL PLAN}
\spectable
\begin{tabularx}{\linewidth}{L{0.24\linewidth}Y}
\toprule
\textbf{Source} & \textbf{House destination} \\
\midrule
Laptops 1–5 & Discrete full-range house inputs 1–5 \\
Laptop 6 & Full-range input 6 if available; otherwise a verified alternate route \\
Control laptop & Network control and projected score; no fixed audience-audio output \\
Documentation & Separate stereo fold-down, not the spatial performance bus \\
\bottomrule
\end{tabularx}
\normalsize

\clearpage
\worktype{Piece 2 · Thursday}{The MacNeoPolitan Trio}{Three laptops, one downbeat, held by the clock rather than a cable.}

\begin{center}
\begin{tikzpicture}[
  every node/.style={font=\sffamily\bfseries\fontsize{9.5pt}{11pt}\selectfont,align=center},
  clock/.style={draw=acpurple,very thick,rounded corners=3pt,fill=acpale,minimum width=60mm,minimum height=12mm},
  mac/.style={very thick,rounded corners=3pt,minimum width=32mm,minimum height=13mm},
  house/.style={draw=acpink,very thick,rounded corners=3pt,fill=acpink!8,minimum width=32mm,minimum height=11mm},
  aud/.style={draw=acdark,very thick,rounded corners=3pt,fill=white,minimum width=44mm,minimum height=11mm},
  flow/.style={-{Stealth[length=3mm]},very thick,draw=acdark}
]
\node[clock] (ntp) at (0,0.8) {outbound NTP · one shared downbeat};
\node[mac,draw=acpurple,fill=acpurple!10] (i) at (-4.5,-1.4) {indigo};
\node[mac,draw=acorange,fill=acorange!10] (c) at (0,-1.4) {citrus};
\node[mac,draw=acpink,fill=acpink!10] (b) at (4.5,-1.4) {blush};
\node[house] (h1) at (-4.5,-3.2) {house input 1};
\node[house] (h2) at (0,-3.2) {house input 2};
\node[house] (h3) at (4.5,-3.2) {house input 3};
\node[aud] (a) at (0,-4.9) {audience in the triangle};
\draw[flow] (ntp) -- (i);
\draw[flow] (ntp) -- (c);
\draw[flow] (ntp) -- (b);
\draw[flow] (i) -- (h1);
\draw[flow] (c) -- (h2);
\draw[flow] (b) -- (h3);
\draw[flow] (h1) -- (a);
\draw[flow] (h2) -- (a);
\draw[flow] (h3) -- (a);
\end{tikzpicture}
\end{center}

Three MacBook Neos — indigo, citrus, and blush — each run Menu Band, the
instrument that lives in the macOS menu bar, and each plays one voice of a
trio. The piece runs from a \texttt{.mbscore}: a JSON score declaring three
machines and one voice each; all three fire at a shared start epoch.

\specbar{SYSTEM REQUIREMENTS}
\spectable
\begin{tabularx}{\linewidth}{L{0.20\linewidth}Y}
\toprule
\textbf{Machines} & Three MacBook Neos running Menu Band on macOS. Artist-supplied, with the score pre-loaded. \\
\textbf{Audio} & \textbf{Three discrete full-range positions} forming a triangle around the
audience, as wide as the room allows, and an interface taking three line-level
feeds from the machines' own outputs: indigo, citrus, and blush to house
inputs 1, 2, and 3. \\
\textbf{Network} & \textbf{Outbound NTP must be reachable.} The machines lock to one downbeat
through synced wall clocks, typically tens of milliseconds across the fleet.
There is no audio-clock link and no timecode between them, so a network that
blocks NTP smears the downbeat. \\
\textbf{Control} & Whether local clients can reach each other over ssh on the house subnet.
Conducting from one host uses ssh; if clients are isolated, the fallback is
arming each machine locally, decided before the technical rehearsal rather than
during it. \\
\textbf{Power} & Three positions. \\
\textbf{Space} & Three stable surfaces at the triangle points. \\
\textbf{Does not use} & USB boot. These machines run macOS from their own disks; AC OS is x86\_64 UEFI only and will not boot on Apple silicon. \\
\bottomrule
\end{tabularx}
\normalsize

\clearpage
\packettitle{The evening}{Running order · changeover · coda}

\pageclaim{The swap happens in front of the audience.}

\begin{center}
\begin{tikzpicture}[
  every node/.style={font=\sffamily\bfseries\fontsize{9pt}{10.5pt}\selectfont,align=center},
  tlstep/.style={draw=acpurple,very thick,rounded corners=3pt,fill=acpale,minimum width=25mm,minimum height=16mm},
  tlswap/.style={draw=acpink,very thick,rounded corners=3pt,fill=acpink!10,minimum width=25mm,minimum height=16mm},
  flow/.style={-{Stealth[length=3mm]},very thick,draw=acdark}
]
\node[tlstep] (p1) at (0,0) {\textbf{Piece 1}\\six laptops};
\node[tlswap] (sw) at (2.9,0) {\textbf{swap}\\5–10 min\\six out, three in};
\node[tlstep] (p2) at (5.8,0) {\textbf{Piece 2}\\three Neos};
\node[tlstep] (talk) at (8.7,0) {conversation};
\node[tlstep] (play) at (11.6,0) {free play\\audience};
\draw[flow] (p1) -- (sw);
\draw[flow] (sw) -- (p2);
\draw[flow] (p2) -- (talk);
\draw[flow] (talk) -- (play);
\end{tikzpicture}
\end{center}

The move between the pieces is a hardware swap, not a patch change: six
machines out, three in. The guest set that used to cover it has been cut, so it
needs five to ten budgeted minutes and a decision about what the room does
meanwhile — house lights and talking, a held drone from the outgoing machines,
or the score on the projector. Please advise what CultureHub prefers.

Six line-level outputs cover both pieces if the patch can be changed during the
swap. If it cannot, nine discrete paths are needed.

After the music there is a conversation, and then the instruments stay on for
the audience to play, so the rig should not be struck until the room clears.

\clearpage
\packettitle{House request}{The combined list}

\section{Audio and network}
\begin{itemize}
\item Kalio 5.1 system with six independently addressable inputs; sixth-channel bandwidth confirmed.
\item Interface or interfaces providing \textbf{six} line-level outputs, with the patch changed during the swap; \textbf{nine} if it cannot be changed.
\item DI boxes, cabling, and adapters appropriate to the confirmed interfaces.
\item Dedicated router or switch isolated from public traffic.
\item \textbf{Outbound NTP reachable}, and confirmation of client-to-client ssh on the house subnet.
\item \textbf{Guest Wi-Fi with internet access and app downloads permitted}, sized for 25–35 devices on the Saturday.
\end{itemize}

\section{Power, room, projection}
\begin{itemize}
\item \textbf{Nine} laptop power positions for the performance, one for the control laptop, conditioned power for audio and network gear, and \textbf{participant power strips for about 20 machines} on the Saturday.
\item Six low tables or plinths around the audience field, plus \textbf{three stable surfaces} for the MacBook Neos.
\item Projector, screen or projection wall, and an HDMI input near the performance position.
\item Central seating for approximately 20–40 people, subject to room capacity.
\item Snack table and bin for the workshop.
\item Stereo livestream and documentation feed derived separately from the spatial mix.
\item One technical rehearsal with CultureHub's audio and streaming staff, long enough to change the patch between the pieces and run both.
\end{itemize}

\section{Jeffrey brings}
\begin{itemize}
\item Six tested surplus laptops and two spares; six AC OS USB drives and two spares, plus workshop loaner sticks.
\item Three MacBook Neos, each with Menu Band installed and the score pre-loaded.
\item Loaner laptops for workshop participants without a machine.
\item Power supplies for all nine performance machines, and the primary control laptop.
\item Owned USB, HDMI, and network adapters; software images and offline recovery media.
\end{itemize}

\clearpage
\packettitle{Open questions}{What CultureHub confirms · what the artist is carrying}

\section{Confirm with CultureHub}
\begin{itemize}
\item Exact Kalio speaker layout, sixth-channel bandwidth, and whether the system accepts six discrete inputs or requires a house spatialization interface.
\item Available interface channel count, connectors, and sample rate, and \textbf{whether the patch can be changed during the swap}.
\item Wired network availability and house network policy — \textbf{outbound NTP, and client-to-client ssh}.
\item \textbf{Whether guest Wi-Fi permits app downloads, and how many devices it holds.}
\item \textbf{Whether participant laptops may be booted from USB in the space.}
\item Projector resolution, throw, connectors, and screen location.
\item Room dimensions, capacity, furniture, and accessible seating route.
\item Livestream platform, capture format, archive ownership, and file delivery.
\item Load-in, soundcheck, public-program, and strike hours, and how much changeover time is available.
\item \textbf{Whether September 19 stays a separate event or folds into September 24.}
\end{itemize}

\section{Risks the artist is carrying}
\begin{itemize}
\item \textbf{The citrus machine is the working development Mac and has 8 GB of RAM}, documented as exhausting memory under load. A hang during Piece 2 stops one of three voices. Mitigation before load-in: perform from a clean boot with nothing else running, or swap in a fourth machine and keep citrus as the spare.
\item \textbf{No spare third Neo yet.} The blush machine is unpurchased at the time of writing; until it exists there is no redundancy for the trio.
\item \textbf{Clock drift is the silent failure.} Unlike a dropped cable, a smeared downbeat still plays and just sounds wrong. It gets verified at the technical rehearsal, on the room's actual network.
\item \textbf{The workshop depends on someone else's network.} Two of the three ways in need working guest Wi-Fi and the third needs USB-boot permission. If all three are blocked the session degrades to the artist's loaner fleet. This is not fixable on the day.
\end{itemize}

\palsfooter
\credit{Aesthetic Computer \& Pals · rider revised 18 August 2026}
\end{document}
